\section{A supporting inequality on five components}\label{sec:certificate}

The two local inequalities use the same central potential. We can
therefore combine their functions into one potential on five disjoint
components. To identify its entire equality set, we must also show
strict inequality for every component pattern other than permutations
of the two opposite-pair triples \((0,a_k,-a_k)\).

\begin{proposition}\label{prop:certificate}
There are pairwise disjoint open balls \(U_0,U_k^+,U_k^-\) centered
at \(0,a_k,-a_k\), respectively, and a bounded continuous function
\(u:U\to\R\), where
\[
 U=U_0\cup U_1^+\cup U_1^-\cup U_2^+\cup U_2^-,
\]
such that
\begin{equation}\label{eq:certificate}
 c(x_1,x_2,x_3)\ge\sum_{i=1}^3u(x_i)
 \quad\text{on }U^3.
\end{equation}
Equality holds precisely at the triples in \(U^3\) which, after a
permutation, have the form
\begin{equation}\label{eq:contact}
 (X_k(y),y,Z_k(y))\in U_0\times U_k^+\times U_k^-,
 \qquad y\in U_k^+,\quad k\in\{1,2\}.
\end{equation}
The maps \(X_k,Z_k\) are diffeomorphisms on \(U_k^+\), with
\[
 X_k(U_k^+)\subset U_0,\qquad Z_k(U_k^+)\subset U_k^-.
\]
\end{proposition}

\begin{proof}
Use the functions and neighborhoods of Lemma~\ref{lem:local}.
Choose small balls contained in these neighborhoods, taking
\(U_0\) inside both central neighborhoods. On the respective
components define
\[
 u=u_0\ \text{on }U_0,\qquad
 u=w_k\ \text{on }U_k^+,\qquad
 u=v_k\ \text{on }U_k^-.
\]
The balls will be shrunk below. Since the functions are smooth near
their centers, we can first restrict to fixed preliminary balls on
which their absolute values are bounded by one constant \(M>0\).
Take the final balls sufficiently small that their closures are
pairwise disjoint.

At the five centers the values of \(u\) are \(5/2,0,0,0,0\).
For a triple of distinct centers other than a permutation of
\((0,a_k,-a_k)\), there are two
possibilities. If zero is absent, its potential sum is zero and
its cost is positive. If zero is present, the other two centers
come from different opposite pairs. Their cost is
\[
 2+\frac1{\sqrt2}>\frac52,
\]
so again there is a strict gap. There are finitely many such
triples, and the cost and local potentials are continuous near
each one. All these strict inequalities therefore persist on
the corresponding products after sufficiently small restrictions.

For a repeated-center type, two coordinates belong to the same
ball. If all radii are at most \(r\), the distance between those
two coordinates is at most \(2r\), and their pair cost is at least
\(1/(2r)\), with the usual interpretation at coincidence.
Taking \(r\) small enough that \(1/(2r)>3M\) makes the cost
strictly larger than the potential sum on every such product.

The inequalities of Lemma~\ref{lem:local} remain valid when their
domains are restricted, with the same equality characterization.
Permuting a triple leaves both \(c\) and the sum of the three
values of \(u\) unchanged. Thus the two prescribed component types and
their permutations satisfy
\eqref{eq:certificate}, with equality exactly as in
\eqref{eq:contact}; every other type has strict inequality.
Finally, keep \(U_0,U_k^-\) fixed and shrink each \(U_k^+\) once
more. Since \(X_k(a_k)=0\) and \(Z_k(a_k)=-a_k\), continuity
ensures
\[
 X_k(U_k^+)\subset U_0,\qquad Z_k(U_k^+)\subset U_k^-.
\]
All inequalities and diffeomorphism properties survive this
restriction. The function \(u\) is bounded and continuous on
the final disjoint union \(U\).
\end{proof}

We now use the central point as parameter for both equality families.
This gives four diffeomorphisms on one common ball, as required by the
selection obstruction and the construction of the smooth marginal.

\begin{lemma}\label{lem:common-ball}
There is an open ball \(B\) of positive radius centered at zero,
contained in \(U_0\), on which the maps
\[
 Y_k=(X_k|_{U_k^+})^{-1}|_B,\qquad W_k=Z_k\circ Y_k
 \quad(k=1,2)
\]
are well defined and are diffeomorphisms onto open images.
They satisfy
\[
 Y_k(B)\subset U_k^+,\qquad W_k(B)\subset U_k^-,
\]
and every triple \((x,Y_k(x),W_k(x))\), \(x\in B\), is an
equality triple in \eqref{eq:certificate}.
In particular, the five sets
\[
 B,\quad Y_1(B),\quad W_1(B),\quad Y_2(B),\quad W_2(B)
\]
are pairwise disjoint.
\end{lemma}

\begin{proof}
Both \(X_k(U_k^+)\) are open neighborhoods of zero. Choose
a positive-radius ball \(B\) about zero inside their intersection
and \(U_0\). The inverse maps then take \(B\) into \(U_k^+\),
and Proposition~\ref{prop:certificate} gives
\(Z_k(Y_k(B))\subset U_k^-\).
The inverse and composition of the relevant diffeomorphisms are
diffeomorphisms onto open images. The identities
\(X_k(Y_k(x))=x\) and \(W_k=Z_k\circ Y_k\) give equality by
Proposition~\ref{prop:certificate}. Disjointness follows from the
five component balls.
\end{proof}
